\section*{Hoofdstuk \ref{ch:g11:seismic-structure} --- Seismische golven onthullen het binnenste}

\begin{solution}{exo:g11:seismic-structure:1}
Een \omterm{def:g11:seismic-structure:body-waves}{P-golf} drukt het gesteente samen en rekt het uit, waarbij het gesteente langs
de voortplantingsrichting beweegt; een \omterm{def:g11:seismic-structure:body-waves}{S-golf} vervormt het door afschuiving, en
het gesteente beweegt loodrecht op die richting. \omterm{def:g11:seismic-structure:body-waves}{P-golven} gaan door vaste stoffen
en vloeistoffen; \omterm{def:g11:seismic-structure:body-waves}{S-golven} alleen door vaste stoffen. De \omterm{def:g11:seismic-structure:body-waves}{P-golf} wordt het eerst
geregistreerd, omdat ze de snelste van de twee is.
\end{solution}

\begin{solution}{exo:g11:seismic-structure:2}
$40 \times 111.2 = \qty{4448}{km}$, ongeveer \qty{4450}{km}. $8900/111.2
\approx \ang{80}$.
\end{solution}

\begin{solution}{exo:g11:seismic-structure:3}
\qty{456}{s} is 7~minuten 36~seconden; \qty{822}{s} is 13~minuten
42~seconden. De S--P-vertraging is $822 - 456 = \qty{366}{s}$, dat is 6~minuten
6~seconden.
\end{solution}

\begin{solution}{exo:g11:seismic-structure:4}
$\sin i_c = 6.0/8.0 = 0.75$, dus $i_c \approx \ang{48.6}$.
\end{solution}

\begin{solution}{exo:g11:seismic-structure:5}
De buitenkern is $5150 - 2891 = \qty{2259}{km}$ dik. De straal van de kern is
$6371 - 2891 = \qty{3480}{km}$, en $3480/6371 \approx 0.55$: iets meer dan de
helft van de straal van de Aarde.
\end{solution}

\begin{solution}{exo:g11:seismic-structure:6}
6~minuten 6~seconden is \qty{366}{s}: de aardbeving ligt \ang{40} ver, dat is
$40 \times 111.2 \approx \qty{4450}{km}$. Eén vertraging geeft alleen de afstand:
het epicentrum kan overal liggen op een cirkel met die straal rond het station.
Er zijn nog twee stations nodig om het ene punt te vinden dat op de juiste afstand
van alle drie ligt.
\end{solution}

\begin{solution}{exo:g11:seismic-structure:7}
$i_c \approx \ang{48.6}$ en $\cos i_c \approx 0.661$. Directe golf: $200/6.0
\approx \qty{33.3}{s}$. \omterm{ex:g11:seismic-structure:sp}{Kopgolf}: $200/8.0 + 2 \times 35 \times 0.661/6.0
\approx 25.0 + 7.7 = \qty{32.7}{s}$. De \omterm{ex:g11:seismic-structure:sp}{kopgolf} komt het eerst aan, ongeveer
\qty{0.6}{s} eerder.
\end{solution}

\begin{solution}{exo:g11:seismic-structure:8}
$\tan i_c \approx 1.134$: de kritische afstand is $2 \times 35 \times 1.134
\approx \qty{79}{km}$. De \omterm{prop:g11:seismic-structure:crossover}{overgangsafstand} is $2 \times 35 \times
\sqrt{14.0/2.0} = 70\sqrt 7 \approx \qty{185}{km}$. Het station van
\cref{exo:g11:seismic-structure:7} ligt op \qty{200}{km}, voorbij de
\omterm{prop:g11:seismic-structure:crossover}{overgangsafstand}: de \omterm{ex:g11:seismic-structure:sp}{kopgolf} komt het eerst aan, zoals gevonden.
\end{solution}

\begin{solution}{exo:g11:seismic-structure:9}
$\sin i_c = 6.5/8.0 = 0.8125$, $i_c \approx \ang{54.3}$ en $\cos i_c \approx
0.583$. \omterm{prop:g11:seismic-structure:head-wave}{Intercepttijd}: $2 \times 6 \times 0.583/6.5 \approx \qty{1.1}{s}$.
\omterm{prop:g11:seismic-structure:crossover}{Overgangsafstand}: $12\sqrt{14.5/1.5} \approx \qty{37}{km}$. Onder een continent
ligt de \omterm{prop:g11:seismic-structure:crossover}{overgangsafstand} rond \qty{165}{km}, dus stations op enkele tientallen
kilometers van elkaar vangen beide rechten op; onder de oceaan komt de \omterm{ex:g11:seismic-structure:sp}{kopgolf} al
voorbij ongeveer \qty{37}{km} het eerst aan, en is de rechte van de directe golf
maar een paar tientallen kilometers lang: de stations moeten een paar kilometer
van elkaar staan om haar te tekenen.
\end{solution}

\begin{solution}{exo:g11:seismic-structure:10}
$\sin r = (8.06/13.7)\sin\ang{30} \approx 0.294$, dus $r \approx \ang{17.1}$:
de straal buigt naar de loodlijn toe, dat wil zeggen naar het middelpunt van de
Aarde. Hij gaat daarom over een lange weg dicht bij het midden door de kern, en
komt ver aan de andere kant weer uit: P-stralen die de kern binnengaan, komen veel
verder van de aardbeving uit dan de stralen die de kern net missen.
\end{solution}

\begin{solution}{exo:g11:seismic-structure:11}
$12\,742/13.7 \approx \qty{930}{s}$, ongeveer 15.5~minuten, terwijl de golf
\qty{1210}{s} nodig heeft, ongeveer 20~minuten: ze is zo'n 280~seconden te laat.
Zelfs met de hoogste snelheid van de mantel zou de golf veel eerder aankomen, dus
het midden van de Aarde moet haar vertragen. Oldham concludeerde dat de Aarde een
centrale kern heeft waarin \omterm{def:g11:seismic-structure:body-waves}{P-golven} trager gaan dan in de mantel.
\end{solution}

\begin{solution}{exo:g11:seismic-structure:12}
(1) Volgens de propositie is $t_i = 2h\cos i_c/v_1$, en $\cos i_c = \sqrt{1 -
v_1^2/v_2^2} = \sqrt{v_2^2 - v_1^2}/v_2$. Dus
\[
  t_i = \frac{2h\sqrt{v_2^2 - v_1^2}}{v_1 v_2},
  \qquad\text{zodat}\qquad
  h = \frac{t_i\,v_1 v_2}{2\sqrt{v_2^2 - v_1^2}}.
\]
(2) $h = 6.23 \times 6.2 \times 8.1/\bigl(2\sqrt{65.61 -
38.44}\bigr) \approx 312.9/10.43 \approx \qty{30.0}{km}$.
\end{solution}

\begin{solution}{exo:g11:seismic-structure:13}
(1) $(3480/6371)^3 \approx 0.163$: de kern is ongeveer \qty{16}{\percent} van
het volume. (2) De gemiddelde dichtheid is het naar volume gewogen gemiddelde:
$5.52 = 0.163
\rho_c + 0.837 \times 4.45$, dus $\rho_c \approx (5.52 - 3.72)/0.163 \approx
\qty{11.0}{g/cm^3}$. (3) Dat is meer dan \qty{9.9}{g/cm^3}, zoals het hoort,
want de dichtheid van de kern neemt toe naar het middelpunt (tot
\qty{13.1}{g/cm^3}); het is ongeveer vier keer de dichtheid van de korst,
\qty{2.6}{g/cm^3}. De kern bestaat uit veel dichter materiaal dan de gesteenten
aan het oppervlak.
\end{solution}

\begin{solution}{exo:g11:seismic-structure:14}
(1) $105 \times 111.2 \approx \qty{11700}{km}$ en $143 \times 111.2 \approx
\qty{15900}{km}$. (2) $(\cos\ang{105} - \cos\ang{143})/2 \approx (-0.259 +
0.799)/2 \approx 0.27$: ongeveer \qty{27}{\percent} van het oppervlak, een ring
rond het punt tegenover de aardbeving. (3) De \omterm{def:g11:seismic-structure:shadow-zone}{schaduwzone} is een bereik van
afstanden tot het epicentrum: ze verschuift met de aardbeving mee. Elke aardbeving
heeft haar eigen ring.
\end{solution}

\begin{solution}{exo:g11:seismic-structure:15}
(1) $h = 9.5 \times 6.3 \times 8.1/\bigl(2\sqrt{65.61 - 39.69}\bigr) \approx
484.8/10.18 \approx \qty{48}{km}$. (2) $\sin i_c = 6.3/8.1 \approx 0.778$,
$\cos i_c \approx 0.629$, en $t_i = 2 \times 30 \times 0.629/6.3 \approx
\qty{6.0}{s}$. (3) De korst onder de keten is ongeveer \qty{48}{km} dik, dikker
dan het continentale gemiddelde van ongeveer \qty{41}{km} en dicht bij het
gemiddelde van gebergtegordels: hoge ketens staan op een dikke korst.
\end{solution}

\begin{solution}{pb:g11:seismic-structure:1}
\textbf{1.} De punten van Pg liggen op een rechte; de punten van Pn liggen op
een andere rechte, minder steil.

\textbf{2.} De aardbeving ligt op heel geringe diepte onder het epicentrum: een
station op afstand nul zou de golf meteen registreren.

\textbf{3.} $(32.3 - 3.3)/(200 - 20) = 29.0/180 \approx \num{0.161}$~\unit{s}
per kilometer.

\textbf{4.} $v_1 = 180/29.0 \approx \qty{6.2}{km/s}$.

\textbf{5.} $3.3 + 280 \times 0.161 \approx \qty{48.4}{s}$.

\textbf{6.} $(43.2 - 18.6)/(300 - 100) = 24.6/200 = \num{0.123}$~\unit{s}
per kilometer, dus $v_2 = 200/24.6 \approx \qty{8.13}{km/s}$.

\textbf{7.} Ze ligt dicht bij \qty{8.0}{km/s}: onder de korst ligt het gesteente
van de bovenste mantel.

\textbf{8.} $\sin i_c = 6.21/8.13 \approx 0.763$, dus $i_c \approx
\ang{49.8}$.

\textbf{9.} De \omterm{ex:g11:seismic-structure:sp}{kopgolf} gaat onder de hoek $i_c$ naar beneden en komt onder
dezelfde hoek weer omhoog. Het eerste punt waar hij naar het oppervlak kan
terugkeren, ligt $h\tan i_c$ voorbij het punt waar hij de \omterm{def:g11:seismic-structure:moho}{Moho} bereikte, dat zelf
op $h\tan i_c$ van de bron ligt: geen station dichter dan $2h\tan i_c$ kan hem
ontvangen.

\textbf{10.} Een seismometer registreert elke golf die hem bereikt, niet alleen
de eerste; de Pn-golf komt na de Pg-golf aan, maar is op de registratie nog als
een afzonderlijke golf te zien.

\textbf{11.} Op \qty{160}{km} komt Pg (\qty{25.8}{s}) vóór Pn
(\qty{26.0}{s}); op \qty{180}{km} komt Pn (\qty{28.5}{s}) vóór Pg
(\qty{29.0}{s}). De wisseling gebeurt tussen deze twee stations.

\textbf{12.} Pg: $t = 3.3 + 0.1611\,(x - 20) \approx 0.08 + 0.1611\,x$. Pn:
$t = 18.6 + 0.123\,(x - 100) = 6.30 + 0.123\,x$ ($x$ in kilometer, $t$ in
seconde).

\textbf{13.} $0.08 + 0.1611\,x = 6.30 + 0.123\,x$ geeft $x \approx
6.22/0.0381 \approx \qty{163}{km}$.

\textbf{14.} Pn, na $6.30 + 0.123 \times 250 \approx \qty{37.1}{s}$.

\textbf{15.} $t_i \approx \qty{6.3}{s}$.

\textbf{16.} $\cos i_c \approx 0.646$, dus $h = 6.21 \times 6.3/(2 \times
0.646) \approx \qty{30.3}{km}$.

\textbf{17.} $x_c = 2 \times 30.3 \times \sqrt{14.34/1.92} \approx 60.6 \times
2.73 \approx \qty{165}{km}$, tot op \qty{2}{km} gelijk aan de \qty{163}{km} van
vraag~13.

\textbf{18.} $\tan i_c \approx 1.18$: $2 \times 30.3 \times 1.18 \approx
\qty{72}{km}$. Het dichtstbijzijnde station dat Pn registreert, ligt op
\qty{80}{km}, en dat op \qty{60}{km} niet: de twee kloppen met elkaar.

\textbf{19.} De tijden zijn tot op \qty{0.1}{s} afgelezen, en een fout van
\qty{0.1}{s} op $t_i$ alleen verandert $h$ al met ongeveer \qty{0.5}{km}; de
snelheden, afgeleid uit hellingen, zijn ook onzeker, en de aardbeving lag niet
precies aan het oppervlak. Bovendien is de korst niet volmaakt uniform. Een
nauwkeurigheid van een of twee kilometer is het hoogste wat de gegevens toelaten.

\textbf{20.} \textbf{De \omterm{def:g11:seismic-structure:moho}{Moho} ligt ongeveer \qty{30}{km} onder dit gebied}: de
korst is hier aan de dunne kant voor een continent, waarvan de korst gemiddeld
ongeveer \qty{41}{km} dik is; \qty{30}{km} is typisch voor gebieden waar de korst
is uitgerekt.
\end{solution}
